\section{Proposed Methodology} \label{sec:methodology}

We propose comprehensive evaluation of three plagiarism detection approaches spanning distinct paradigms: classical machine learning, deep neural networks, and pre-trained transformers. Each approach represents a different point on the accuracy-efficiency frontier.

\subsection{TF-IDF Baseline Approach}

\subsubsection{Theoretical Foundation}

TF-IDF (Term Frequency-Inverse Document Frequency) quantifies text similarity through word importance weighting. For document $D_i$, TF-IDF is computed as:

\begin{equation}
\text{TF-IDF}(t, D_i) = \text{TF}(t, D_i) \times \log\left(\frac{N}{|\\{D_j : t \in D_j\\}|}\right)
\end{equation}

where $N$ is total documents, and $|\\{D_j : t \in D_j\\}|$ is document frequency of term $t$. Document similarity uses cosine distance:

\begin{equation}
\text{similarity}(D_a, D_b) = \frac{\vec{v}_a \cdot \vec{v}_b}{|\vec{v}_a| \cdot |\vec{v}_b|}
\end{equation}

\subsubsection{Implementation Details}

\begin{itemize}
  \item \textbf{Vectorization} --- scikit-learn TfidfVectorizer with max\_features=5000
  \item \textbf{N-gram range} --- (1,2) to capture bigrams alongside unigrams
  \item \textbf{Decision threshold} --- 0.5 (similarity $\geq 0.5$ indicates plagiarism)
  \item \textbf{Complexity} --- $O(n \times m)$ where $n$ is document count, $m$ is feature dimension
\end{itemize}

\subsection{BiLSTM Deep Learning Approach}

\subsubsection{Architecture Design}

BiLSTM (Bidirectional Long Short-Term Memory) captures sequential dependencies bidirectionally:

\begin{equation}
\vec{h}_t = \text{BiLSTM}(x_t, h_{t-1}) = [\overrightarrow{h}_t \oplus \overleftarrow{h}_t]
\end{equation}

Our architecture:
\begin{enumerate}
  \item \textbf{Embedding layer} --- 100-dimensional embeddings for vocabulary
  \item \textbf{BiLSTM layer} --- 64 hidden units with dropout=0.3
  \item \textbf{Pooling} --- Global max pooling across temporal dimension
  \item \textbf{Classification head} --- 3 fully-connected layers (128 → 64 → 32 → 1)
  \item \textbf{Activation} --- ReLU for hidden layers, Sigmoid for output
\end{enumerate}

\subsubsection{Training Configuration}

\begin{itemize}
  \item \textbf{Optimizer} --- Adam with learning rate 0.001
  \item \textbf{Loss function} --- Binary Cross-Entropy
  \item \textbf{Epochs} --- 10 with early stopping (patience=2)
  \item \textbf{Batch size} --- 32 samples per gradient update
  \item \textbf{Regularization} --- L2 weight decay (0.0001), Dropout (0.3)
  \item \textbf{Framework} --- PyTorch with CUDA GPU acceleration
\end{itemize}

\subsection{BERT Transformer Approach}

\subsubsection{Pre-trained Model Selection}

We employ bert-base-uncased (110M parameters), pre-trained on English Wikipedia and BookCorpus. This 12-layer transformer model captures bidirectional context through multi-head self-attention mechanisms.

\subsubsection{Fine-tuning Strategy}

For task-specific plagiarism detection:

\begin{enumerate}
  \item \textbf{Task head} --- Add 2-layer MLP classification head atop [CLS] token representation
  \item \textbf{Fine-tuning} --- Update all BERT parameters and task-specific head jointly
  \item \textbf{Input format} --- Concatenate document pairs: ``[CLS] document1 [SEP] document2 [SEP]''
  \item \textbf{Maximum length} --- 128 tokens (truncate longer documents)
\end{enumerate}

\subsubsection{Training Configuration}

\begin{itemize}
  \item \textbf{Learning rate} --- 2e-5 (lower than BiLSTM to preserve pre-training)
  \item \textbf{Epochs} --- 5 with validation-based early stopping
  \item \textbf{Batch size} --- 16 (larger batches cause OOM errors)
  \item \textbf{Optimizer} --- AdamW with weight decay
  \item \textbf{Warmup steps} --- 10\% of total steps for learning rate scheduling
  \item \textbf{Framework} --- Hugging Face Transformers library with PyTorch backend
\end{itemize}

\subsection{Shared Evaluation Metrics}

All models are evaluated on five complementary metrics:

\begin{enumerate}
  \item \textbf{Accuracy} --- Overall correctness: $\frac{TP + TN}{TP + TN + FP + FN}$
  
  \item \textbf{Precision} --- False positive rate control: $\frac{TP}{TP + FP}$
  
  \item \textbf{Recall} --- Detection completeness: $\frac{TP}{TP + FN}$
  
  \item \textbf{F1-Score} --- Harmonic mean: $2 \times \frac{\text{Precision} \times \text{Recall}}{\text{Precision} + \text{Recall}}$
  
  \item \textbf{ROC-AUC} --- Trade-off analysis: Area under receiver operating characteristic curve
\end{enumerate}

All metrics computed on held-out test set to ensure unbiased evaluation.
